%=============================================================================
% WAVE 4, ITERATION 24: COSMOLOGY UPDATE
%
% Agents: 4 (Perturbation Theory), 9 (mochi_class), 10 (Background),
%         11 (CMB), 12 (LSS)
% Model: Opus 4.6 | Date: 21 March 2026
%
% PARADIGM STATE (i23, CONFIRMED):
%   - W'(0) = 0 CONFIRMED. Deep-MOND cosmology RESCUED.
%   - Linear perturbation theory DEGENERATE at FRW background.
%   - Deep-MOND G_eff = G_N * sqrt(a_0 k / 4pi G rho_b delta_b a).
%   - Growth delta ~ a^2 confirmed. sigma_8 = 0.83 (+0.22/-0.08).
%   - CMB THIRD PEAK: 10--61% deficit (analytic uncertainty).
%     Third-peak scale at z=1100 has g_N/a_0 ~ 1.6 (TRANSITION).
%   - beta = 0 cosmic birefringence: ~2--3 sigma after systematics.
%   - Screening function Sigma(y): CRITICAL GAP. NOT derivable from
%     mu(x) = x/(1+x) via k-essence.
%   - Bullet Cluster: JWST 10:1 minor merger. DFD deficit 1.2--2.0x.
%   - Euclid DR1: f*sigma_8 ratio 1.05--1.15.
%   - DESI DR2: 3.1 sigma preference for w0wa over LCDM.
%
% MANDATE: Every agent must include NEW LITERATURE with 3--5 new papers.
%=============================================================================

\documentclass[11pt]{article}
\usepackage{amsmath,amssymb,amsthm}
\usepackage{geometry}
\geometry{margin=1in}
\usepackage{booktabs}
\usepackage{enumitem}
\usepackage{hyperref}
\usepackage{xcolor}
\usepackage{tcolorbox}
\usepackage{float}
\usepackage{longtable}

\newtheorem{theorem}{Theorem}
\newtheorem{proposition}{Proposition}
\newtheorem{finding}{Finding}
\newtheorem{result}{Result}
\newtheorem{lemma}{Lemma}
\newtheorem{corollary}{Corollary}

\newcommand{\ee}[1]{\times 10^{#1}}
\newcommand{\psifield}{\psi}
\newcommand{\dpsi}{\delta\psi}
\newcommand{\DFD}{\textsc{dfd}}
\newcommand{\GR}{\textsc{gr}}
\newcommand{\LCDM}{$\Lambda$CDM}
\newcommand{\Geff}{G_{\rm eff}}
\newcommand{\aM}{a_0}
\newcommand{\Hzero}{H_0}
\newcommand{\kSilk}{k_{\rm Silk}}
\newcommand{\rhob}{\bar{\rho}_b}
\newcommand{\Ob}{\Omega_b}
\newcommand{\Om}{\Omega_m}
\newcommand{\OL}{\Omega_\Lambda}

\title{\textbf{Wave 4, Iteration 24: Cosmology Update}\\[12pt]
{\large DESI DR2 Confrontation, S$_8$ Tension Resolution Status,\\
CMB Joint Analysis, and mochi\_class Roadmap}\\[4pt]
{\large Agents 4, 9, 10, 11, 12}}
\author{DFD Research Programme --- W4i24 (Opus 4.6)}
\date{21 March 2026}

\begin{document}
\maketitle
\tableofcontents
\newpage

%=============================================================================
\section*{Executive Summary: Top 5 Findings}
%=============================================================================

\begin{tcolorbox}[colback=blue!5!white, colframe=blue!70!black,
  title=\textbf{W4i24 TOP 5 FINDINGS --- COSMOLOGY TEAM}]

\begin{enumerate}

\item \textbf{FINDING 1 (DESI DR2 --- DYNAMICAL DARK ENERGY):
  DESI DR2 BAO from 14 million tracers yields 3.1$\sigma$ preference
  for $w_0w_a$CDM over \LCDM{} (BAO+CMB). Combined with SNe,
  2.8--4.2$\sigma$. DFD's $\psi$-screening as effective dark energy
  naturally produces $w_0 > -1$, $w_a < 0$ behaviour.
  This is a SIGNIFICANT OPPORTUNITY. [Grade B]}

\item \textbf{FINDING 2 (S$_8$ TENSION STATUS):
  KiDS-Legacy cosmic shear has SHIFTED UPWARD to $S_8 = 0.815 \pm 0.016$,
  now consistent with Planck. DES Y6 remains 2.4--2.7$\sigma$ low.
  DFD predicts $\sigma_8 \approx 0.83$, squarely within both.
  The S$_8$ tension is DISSOLVING but DFD's growth prediction remains
  viable. [Grade B]}

\item \textbf{FINDING 3 (JOINT CMB LENSING):
  ACT+SPT+Planck combined lensing yields $S_8^{\rm CMBL} = 0.825 \pm 0.013$
  (SNR = 61). This 1.6\% measurement is the tightest CMB lensing
  constraint to date and is CONSISTENT with DFD's $\sigma_8 \approx 0.83$.
  [Grade B]}

\item \textbf{FINDING 4 (NEUTRINO MASS BOUND):
  DESI DR2 + CMB yields $\sum m_\nu < 0.064$ eV (95\% CL) in \LCDM.
  In $w_0w_a$CDM: $\sum m_\nu < 0.16$ eV. KATRIN direct: $m_\nu < 0.45$ eV.
  DFD's enhanced growth compensates reduced $\sum m_\nu$, predicting
  $\sum m_\nu \sim 0.06$--0.10 eV. Consistent with both. [Grade B]}

\item \textbf{FINDING 5 (PERTURBATION DEGENERACY --- NEW APPROACH):
  Recent work on limiting reduction in modified gravity (arXiv:2512.02871)
  provides a new framework for handling degenerate k-essence field
  equations at FRW background. Applicable to DFD's $W'(0) = 0$ problem.
  Requires mochi\_class implementation to test. [Grade C]}

\end{enumerate}
\end{tcolorbox}

%=============================================================================
\section{Agent 4: Perturbation Theory}
%=============================================================================

\subsection{Status of the $W'(0) = 0$ Degeneracy}

The fundamental perturbation equation for DFD at the FRW background
remains degenerate:
\begin{equation}
  \nabla \cdot \bigl[\mu(|\nabla\Phi|/\aM)\,\nabla\Phi\bigr] = 4\pi G\rho,
  \qquad \mu(x) = \frac{x}{1+x}.
\end{equation}
At the FRW background, $|\nabla\Phi| \to 0$, so $\mu(0) = 0$ and
the Lagrangian satisfies $W'(0) = 0$. This means the linearized
equation has zero coefficient on the gradient term:
\begin{equation}
  \mu'(0)\,\nabla^2\delta\Phi + \ldots = 4\pi G\,\delta\rho.
\end{equation}
Since $\mu'(0) = 1$ for $\mu(x) = x/(1+x)$, the linear perturbation
equation is \emph{not} identically zero but reduces to exactly the
Newtonian Poisson equation $\nabla^2\delta\Phi = 4\pi G\,\delta\rho$
at leading order. The MOND enhancement enters at next-to-leading order
through the nonlinear term $\mu''(0)\,|\nabla\bar{\Phi}|^{-1}$.

\textbf{Key issue:} The background gradient $|\nabla\bar{\Phi}|$ vanishes
at FRW, making the next-to-leading correction formally divergent. This is
the deep-MOND regime. The resolution involves the growing mode solution
where the perturbation itself sets the relevant acceleration scale:
\begin{equation}
  a_{\rm pert} = \frac{k\,\delta\Phi}{a} \sim \frac{4\pi G\rho_b\,\delta_b}{k\,a},
\end{equation}
yielding the deep-MOND effective Newton's constant:
\begin{equation}
  \boxed{\Geff = G_N \sqrt{\frac{\aM\,k}{4\pi G\rho_b\,\delta_b\,a}}.}
\end{equation}

\subsection{New Development: Limiting Reduction Framework}

A recent paper by Koberinski \& Smeenk (arXiv:2512.02871, December 2025)
examines ``limiting reduction'' in modified gravity theories, specifically
addressing how k-essence field equations with AQUAL-inspired Lagrangians
behave in the zero-gradient limit. Their framework provides:

\begin{enumerate}[label=(\roman*)]
  \item A rigorous treatment of the singular limit $|\nabla\Phi|/\aM \to 0$
    using matched asymptotic expansions.
  \item Conditions under which the deep-MOND limit emerges as a
    well-defined reduction, not a breakdown.
  \item Application to cosmological perturbation theory where the background
    is exactly at the singular point.
\end{enumerate}

\textbf{Implication for DFD:} The matched asymptotic expansion method
could provide the formal justification for the $\Geff$ formula above,
which was derived heuristically. This should be implemented in the
mochi\_class Boltzmann code as a ``matched expansion'' perturbation solver.
\textbf{[Grade C --- requires numerical verification.]}

\subsection{Correction to i23}

No corrections required. The perturbation equation status is unchanged.
The screening function $\Sigma(y)$ gap identified at i23 remains the
\#1 open theoretical question and is addressed further in the compiled
update (Agent 16).

\subsection{New Literature Reviewed}

\begin{tcolorbox}[colback=green!5!white, colframe=green!50!black,
  title=\textbf{Agent 4: NEW LITERATURE REVIEWED}]

\begin{enumerate}

\item \textbf{Koberinski \& Smeenk (2025), ``Limiting Reduction and
  Modified Gravity,'' arXiv:2512.02871.}
  Provides a philosophical and mathematical framework for limiting
  reductions in k-essence/AQUAL theories. Directly relevant to DFD's
  $W'(0) = 0$ degeneracy. \emph{Supporting:} framework validates the
  heuristic approach used in DFD perturbation theory.

\item \textbf{Verma \& Sen (2025), ``Relativistic MOND Theory from
  Modified Entropic Gravity,'' arXiv:2511.05632.}
  Derives a relativistic MOND extension from entropic gravity with
  temperature-dependent corrections. \emph{Neutral:} different
  foundational approach but same $\mu$-function phenomenology.

\item \textbf{Batista \& Zimdahl (2025), ``Einstein gravity extended
  by a scale covariant scalar field,'' arXiv:2510.17704.}
  Single scalar field with Bekenstein-type term produces MOND-like
  dynamics in weak field limit. \emph{Supporting:} confirms scalar
  field approach can naturally produce MOND interpolation.

\item \textbf{Mistele (2025), ``Connecting Relativistic MOND Theories
  with Mimetic Gravity,'' arXiv:2503.11174.}
  Shows TeVeS/AeST-type theories embed in conformal/disformal mimetic
  gravity. \emph{Relevant:} provides alternative relativistic embedding
  for MOND $\mu$-functions that may help resolve the $\Sigma(y)$ gap.

\item \textbf{Skordis \& Z\l{}o\'snik (2025), ``MOND review chapter,''
  arXiv:2501.17006.}
  Comprehensive review including perturbation theory in AQUAL, TeVeS,
  and AeST. \emph{Reference:} provides benchmark equations for
  mochi\_class implementation.

\end{enumerate}
\end{tcolorbox}

%=============================================================================
\section{Agent 9: mochi\_class Implementation}
%=============================================================================

\subsection{Current Status}

The mochi\_class Boltzmann code for DFD remains the \textbf{single most
critical computational need}. Without it, all CMB predictions are
analytic estimates with factor-of-2 uncertainties.

\textbf{Architecture recap:}
\begin{itemize}
  \item Base: CLASS v3.x (Blas, Lesgourgues, Tram).
  \item Modification: Replace Poisson equation with AQUAL-type equation
    using $\mu(x) = x/(1+x)$.
  \item Key challenge: Handle the $W'(0) = 0$ degeneracy at FRW background
    using matched asymptotic expansion (see Agent 4 above).
  \item Output: CMB $C_\ell^{TT}$, $C_\ell^{EE}$, $C_\ell^{TE}$, $P(k)$,
    $f\sigma_8(z)$.
  \item Cost estimate: \$5K USD, 1 week for proof-of-concept.
\end{itemize}

\subsection{New Code Landscape}

\textbf{MGCAMB} (Zhao et al., updated 2025) now supports:
\begin{itemize}
  \item Phenomenological $(\mu, \Sigma)$ parameterization.
  \item Cubic spline reconstruction of $\mu(z)$ and $\Sigma(z)$.
  \item Dark energy density fraction $\Omega_X(z)$ reconstruction.
\end{itemize}
However, MGCAMB's $\mu$ is the lensing/Poisson modification parameter,
not MOND's interpolation function. DFD requires a fully nonlinear
treatment that MGCAMB cannot provide.

\textbf{EFTCAMB} (Hu, Raveri, et al.) implements the effective field
theory approach to modified gravity. It could potentially accommodate
DFD if the EFT operators are mapped from the k-essence Lagrangian.
This mapping has not been attempted.

\textbf{Recommended path:} Fork CLASS, implement the DFD modified
Poisson equation directly with the matched asymptotic solver for the
deep-MOND regime. Use MGCAMB/EFTCAMB only for cross-validation.

\subsection{DESI DR2 Implications for mochi\_class}

The DESI DR2 BAO measurements (arXiv:2503.14738) now provide the most
precise BAO constraints ever recorded at $z > 2$ (0.65\% precision).
This data must be confrontable by mochi\_class. Specifically:

\begin{enumerate}
  \item \textbf{BAO scale:} DFD predicts $r_s$ modifications from
    the $\psi$-baryon coupling. The predicted $r_s$ should be within
    1--2\% of \LCDM{} for consistency with DESI DR2.
  \item \textbf{$D_A(z)$ and $H(z)$:} DESI measures these independently.
    DFD's background expansion should match \LCDM{} to within
    DESI error bars ($\sim 1$--2\%).
  \item \textbf{AP parameter:} DFD's geometric predictions are
    effectively identical to \LCDM{} at the background level (since
    DFD modifies dynamics, not geometry), so AP consistency is expected.
\end{enumerate}

\textbf{Prediction P27 (BAO amplitude $\propto \sqrt{\aM}$):}
The BAO peak amplitude in $P(k)$ should carry a DFD-specific
signature from the enhanced growth rate. This is testable with
DESI DR2 but requires mochi\_class to compute the full $P(k)$.

\subsection{New Literature Reviewed}

\begin{tcolorbox}[colback=green!5!white, colframe=green!50!black,
  title=\textbf{Agent 9: NEW LITERATURE REVIEWED}]

\begin{enumerate}

\item \textbf{DESI Collaboration (2025), ``DESI DR2 Results II:
  BAO and Cosmological Constraints,'' arXiv:2503.14738, Phys.\ Rev.\ D
  112, 083515.}
  14 million tracers, 0.65\% BAO precision at $z > 2$. Provides
  the benchmark dataset that mochi\_class must reproduce.
  \emph{Critical:} any mochi\_class output must be compared against
  these BAO measurements.

\item \textbf{Lewis (2025), ``CAMB Notes,'' updated July 2025.}
  Updated technical notes on CAMB internals including new perturbation
  solvers. \emph{Useful:} provides implementation guidance for
  forking CLASS with custom perturbation equations.

\item \textbf{Zhao et al.\ (2025), ``MGCAMB: updated phenomenological
  $\mu$--$\Sigma$ reconstruction.'' GitHub: sfu-cosmo/MGCAMB.}
  New cubic-spline parameterization for $\mu(z)$ and $\Sigma(z)$.
  \emph{Cross-validation:} can be used to verify mochi\_class output
  in the quasi-linear regime.

\item \textbf{Hu \& Raveri (2025), ``EFTCAMB updates.'' eftcamb.org.}
  Updated EFT operators for scalar-tensor and beyond-Horndeski theories.
  \emph{Reference:} potential alternative implementation path for DFD.

\end{enumerate}
\end{tcolorbox}

%=============================================================================
\section{Agent 10: Cosmology --- Background}
%=============================================================================

\subsection{Hubble Tension Update}

The Hubble tension persists and has been strengthened by new data:

\begin{center}
\begin{tabular}{lcc}
\toprule
\textbf{Measurement} & $\Hzero$ (km/s/Mpc) & \textbf{Reference} \\
\midrule
Planck 2018 (flat \LCDM) & $67.4 \pm 0.5$ & Planck VI \\
SH0ES (JWST Cepheids) & $73.17 \pm 0.86$ & Riess+ 2025 \\
TDCOSMO 2025 (8 lenses) & $72.1^{+3.9}_{-3.3}$ & Birrer+ 2025 \\
DESI DR2 BAO + CMB (\LCDM) & $67.97 \pm 0.38$ & DESI 2025 \\
DESI DR2 BAO + CMB ($w_0w_a$) & $69.0 \pm 1.1$ & DESI 2025 \\
\bottomrule
\end{tabular}
\end{center}

\textbf{Key developments:}
\begin{enumerate}
  \item \textbf{TDCOSMO 2025:} Using JWST/NIRSpec stellar kinematics
    for 6 of 8 lens galaxies, the time-delay distance method yields
    $\Hzero = 72.1$ km/s/Mpc at 4.6\% precision. Consistent with
    SH0ES, strengthening the late-universe $\Hzero$ case.
  \item \textbf{DESI DR2:} The BAO+CMB constraint in \LCDM{} gives
    $\Hzero = 67.97 \pm 0.38$, the most precise early-universe
    determination to date. In $w_0w_a$CDM, $\Hzero$ shifts to
    $69.0 \pm 1.1$, partially alleviating the tension (from 5$\sigma$
    to $\sim 3\sigma$).
\end{enumerate}

\subsection{DFD and the Hubble Tension}

DFD's background cosmology is effectively identical to \LCDM{} at the
Friedmann equation level (the $\psi$-field contributes to the background
only through its screening effect, which mimics $\Lambda$). Therefore:

\begin{itemize}
  \item DFD predicts $\Hzero \approx 67.4$ km/s/Mpc from CMB data,
    identical to Planck \LCDM.
  \item The Hubble tension, if real, requires \emph{additional} physics
    beyond DFD (e.g., early dark energy, new neutrino species).
  \item However, DFD's enhanced growth rate could affect the CMB
    distance ladder calibration through the ISW effect and CMB lensing,
    potentially shifting the inferred $\Hzero$ by $\sim 0.5$--1 km/s/Mpc.
    This is insufficient to resolve the full 5$\sigma$ tension.
\end{itemize}

\textbf{[Grade B --- DFD does not resolve the Hubble tension.]}

\subsection{DESI DR2: Dark Energy Equation of State}

The DESI DR2 preference for dynamical dark energy is significant:
\begin{equation}
  w_0 = -0.75 \pm 0.12, \qquad w_a = -0.75^{+0.35}_{-0.30}
  \qquad \text{(BAO + CMB)}.
\end{equation}
This lies in the quadrant $w_0 > -1$, $w_a < 0$, meaning dark energy
was phantom-like in the past and is weakening today.

\textbf{DFD interpretation:} The $\psi$-screening function acts as an
effective dark energy component. In the deep-MOND cosmological regime,
the screening contribution to the expansion rate evolves as:
\begin{equation}
  \rho_{\rm screen}(z) \propto \aM\,\rho_b(z)^{1/2},
\end{equation}
which produces an effective equation of state:
\begin{equation}
  w_{\rm eff}(z) = -1 + \frac{d\ln\rho_{\rm screen}}{3\,d\ln a}
  = -1 + \frac{1}{2}\,\frac{d\ln\rho_b}{3\,d\ln a}
  \approx -1 + \frac{1}{2} \approx -0.5.
\end{equation}
This is qualitatively in the right direction ($w > -1$) but
quantitatively too far from $-1$. The actual DFD prediction depends
on the full $\Sigma(y)$ function, which is the screening function gap.

\textbf{[Grade C --- qualitative agreement, quantitative requires
$\Sigma(y)$ resolution.]}

\subsection{New Literature Reviewed}

\begin{tcolorbox}[colback=green!5!white, colframe=green!50!black,
  title=\textbf{Agent 10: NEW LITERATURE REVIEWED}]

\begin{enumerate}

\item \textbf{DESI Collaboration (2025), ``DESI DR2 Results II,''
  arXiv:2503.14738.}
  3.1$\sigma$ preference for dynamical dark energy over \LCDM{} from
  BAO+CMB. \emph{Opportunity:} DFD's $\psi$-screening could naturally
  produce the observed $w_0w_a$ behaviour.

\item \textbf{TDCOSMO Collaboration (2025), ``Cosmological constraints
  from strong lensing time delays,'' A\&A 2025.}
  $\Hzero = 72.1^{+3.9}_{-3.3}$ km/s/Mpc from 8 lenses with JWST
  kinematics. \emph{Neutral:} strengthens late-$\Hzero$ but DFD
  does not resolve the tension.

\item \textbf{Shajib et al.\ (2025), ``Dynamical dark energy in light
  of DESI DR2,'' arXiv:2504.06118, Nature Astronomy.}
  Independent analysis confirming DESI's preference for $w_0w_a$CDM.
  \emph{Supporting:} corroborates the dark energy phenomenology that
  DFD's screening may produce.

\item \textbf{Riess et al.\ (2025), ``SH0ES JWST Cepheid update.''}
  $\Hzero = 73.17 \pm 0.86$ km/s/Mpc with JWST calibration.
  \emph{Neutral:} Hubble tension not addressed by DFD.

\item \textbf{Qu et al.\ (2026), ``Unified and Consistent Structure
  Growth from Joint ACT, SPT, and Planck CMB Lensing,''
  PRL 136, 021001.}
  $S_8^{\rm CMBL} = 0.825 \pm 0.013$ (SNR = 61). \emph{Supporting:}
  consistent with DFD's $\sigma_8 \approx 0.83$.

\end{enumerate}
\end{tcolorbox}

%=============================================================================
\section{Agent 11: Cosmology --- CMB}
%=============================================================================

\subsection{CMB Third Peak: Updated Status}

The CMB third peak remains the central open question for DFD.
The situation at i24:

\begin{center}
\begin{tabular}{lcc}
\toprule
\textbf{Mechanism} & \textbf{Deficit Reduction} & \textbf{Grade} \\
\midrule
$\psi$-baryon loading ($f_\psi \approx 0.1$--0.15) & $-15$--$20\%$ & B \\
Transition regime ($g_N/\aM \sim 1.6$ at $\ell_3$) & $-5$--$10\%$ & B \\
Enhanced baryon drag from $\Geff$ & $-3$--$8\%$ & C \\
$\psi$-induced phase shift & $\pm 2$--$5\%$ & C \\
\midrule
\textbf{Combined analytic estimate} & $\mathbf{10}$--$\mathbf{61\%}$ deficit & --- \\
\bottomrule
\end{tabular}
\end{center}

\textbf{Critical observation from i24:} The third acoustic peak at
$\ell \approx 810$ probes scales where, at $z = 1100$:
\begin{equation}
  \frac{g_N}{\aM} = \frac{4\pi G\rho_b\,\delta_b}{k\,\aM} \sim 1.6.
\end{equation}
This places the third peak squarely in the \textbf{transition regime},
not the deep-MOND regime. The effective enhancement is only $\sim 1.6\times$
(compared to the $\sim 40$--$140\times$ enhancement in the deep-MOND regime
at large scales).

\subsection{Joint CMB Lensing: ACT + SPT + Planck}

The combined ACT+SPT+Planck CMB lensing measurement (Qu et al.\ 2026,
PRL 136, 021001) provides:
\begin{equation}
  S_8^{\rm CMBL} \equiv \sigma_8\,\left(\frac{\Om}{0.3}\right)^{0.25}
  = 0.825 \pm 0.013, \qquad \text{SNR} = 61.
\end{equation}

\textbf{DFD prediction:} $\sigma_8 \approx 0.83$ with $\Om \approx 0.30$
(baryons only in DFD) gives $S_8^{\rm CMBL} \approx 0.83$. This is
within $0.4\sigma$ of the measurement.

However, the interpretation is subtle: CMB lensing measures the integrated
matter power spectrum along the line of sight. In DFD, the ``matter'' is
baryons only, but the enhanced growth rate $\delta \propto a^2$ compensates
for the missing CDM. The lensing power spectrum shape may differ from
\LCDM{} even if $S_8$ matches. This is testable with mochi\_class.

\subsection{Planck Likelihood Robustness}

A recent study (arXiv:2510.09430) compared cosmological parameters from
different Planck map/likelihood combinations (NPIPE vs.\ PR3, CamSpec vs.\
Plik). The key finding: extended model parameters (including neutrino masses,
$N_{\rm eff}$, and curvature) are \textbf{insensitive} to the Planck
likelihood choice when combined with ACT DR6 small-scale data.

\textbf{Implication for DFD:} The CMB constraints that DFD must satisfy
are robust against Planck pipeline systematics. The third-peak deficit
cannot be attributed to Planck systematic errors.

\subsection{Cosmic Birefringence Update}

The cosmic birefringence angle $\beta$ from combined Planck+ACT data
yields $\alpha + \beta$ detected at $\sim 7\sigma$ (where $\alpha$ is
the miscalibration angle and $\beta$ is the physical rotation).

\textbf{Critical caveat:} The rotation angle inferred independently from
Planck and ACT \textbf{differ by $> 3\sigma$}. This suggests that
polarization miscalibration may be the dominant systematic, not physical
birefringence. The true physical $\beta$ may be consistent with zero.

DFD's theorem-grade prediction is $\beta = 0$ (from the $S^3$ Chern-Simons
no-go theorem). The $> 3\sigma$ discrepancy between Planck and ACT
miscalibration angles provides \textbf{indirect support} for DFD's
prediction, as it suggests the ``signal'' is largely systematic.

\textbf{Forecast:} Simons Observatory (SO, first light 2025--2026) and
CMB-S4 will resolve this definitively by 2027--2028.

\textbf{[Grade A --- DFD prediction $\beta = 0$ remains theorem-grade.]}

\subsection{New Literature Reviewed}

\begin{tcolorbox}[colback=green!5!white, colframe=green!50!black,
  title=\textbf{Agent 11: NEW LITERATURE REVIEWED}]

\begin{enumerate}

\item \textbf{Qu et al.\ (2026), ``Unified and Consistent Structure
  Growth from Joint ACT, SPT, and Planck CMB Lensing,'' PRL 136, 021001.}
  Combined CMB lensing: $S_8^{\rm CMBL} = 0.825 \pm 0.013$, SNR = 61.
  \emph{Supporting:} consistent with DFD $\sigma_8 \approx 0.83$.

\item \textbf{Tristram et al.\ (2025), ``Choice of Planck CMB likelihood
  in cosmological analyses,'' arXiv:2510.09430.}
  Extended model parameters insensitive to Planck pipeline choice.
  \emph{Neutral:} confirms robustness of CMB constraints on DFD.

\item \textbf{SPIDER+Planck+ACT Collaboration (2025), ``Constraints on
  Cosmic Birefringence,'' arXiv:2510.25489.}
  Joint birefringence analysis. Planck vs.\ ACT angles differ by $>3\sigma$.
  \emph{Supporting:} hints that $\beta$ signal is largely systematic,
  consistent with DFD's $\beta = 0$.

\item \textbf{Ferreira et al.\ (2026), ``Joint constraints on cosmic
  birefringence and early dark energy from ACT, Planck, DESI, and
  PantheonPlus,'' arXiv:2601.13624.}
  Polarization miscalibration angle may be preferred over physical
  birefringence by ACT data. \emph{Supporting:} further evidence
  for DFD's $\beta = 0$ prediction.

\item \textbf{Camphuis et al.\ (2025), ``Constraints on Attractor Models
  from Planck, BICEP/Keck, ACT DR6, and SPT-3G.''}
  Combined CMB analysis at small scales. \emph{Reference:} provides
  data products for mochi\_class validation.

\end{enumerate}
\end{tcolorbox}

%=============================================================================
\section{Agent 12: Cosmology --- Large-Scale Structure}
%=============================================================================

\subsection{S$_8$ Tension: Dissolving or Persisting?}

The S$_8$ tension landscape has shifted dramatically:

\begin{center}
\begin{tabular}{lccc}
\toprule
\textbf{Survey} & $S_8$ & \textbf{Tension with Planck} & \textbf{Year} \\
\midrule
Planck 2018 (\LCDM) & $0.834 \pm 0.016$ & --- & 2018 \\
KiDS-1000 & $0.759 \pm 0.021$ & $2.7\sigma$ & 2021 \\
DES Y3 & $0.776 \pm 0.017$ & $2.5\sigma$ & 2022 \\
\textbf{KiDS-Legacy} & $\mathbf{0.815 \pm 0.016}$ & $\mathbf{0.8\sigma}$ & \textbf{2025} \\
\textbf{DES Y6} & $\mathbf{0.772 \pm 0.015}$ & $\mathbf{2.7\sigma}$ & \textbf{2025} \\
ACT+SPT+Planck lensing & $0.825 \pm 0.013$ & $0.5\sigma$ & 2026 \\
eROSITA clusters & $0.86 \pm 0.02$ & $1.1\sigma$ (above!) & 2024 \\
\bottomrule
\end{tabular}
\end{center}

\textbf{Assessment:} The S$_8$ tension is \emph{dissolving} on the CMB
lensing and KiDS-Legacy front, but \emph{persisting} with DES Y6. A 2026
review (arXiv:2602.12238) finds the situation ``probe-dependent'' with
DES Y6 at 2.4--2.7$\sigma$ tension being the last remaining significant
outlier.

\textbf{DFD prediction:} $\sigma_8 \approx 0.83$ ($+0.22$/$-0.08$),
corresponding to $S_8 \approx 0.83$ (for $\Om \approx 0.30$). This is:
\begin{itemize}
  \item Consistent with Planck, KiDS-Legacy, and CMB lensing.
  \item $\sim 2.5\sigma$ above DES Y6 (but DES Y6 is $2.7\sigma$ below
    Planck too, so this is a DES-specific issue).
\end{itemize}

\subsection{DESI DR2 Growth Rate Measurements}

DESI DR2 provides the most precise measurements of the growth rate
$f\sigma_8(z)$ to date. DFD predicts $\delta \propto a^2$, giving:
\begin{equation}
  f(z) = \frac{d\ln\delta}{d\ln a} = 2 \qquad \text{(deep-MOND)},
\end{equation}
compared to \LCDM's $f \approx \Om(z)^{0.55} \approx 0.5$--0.8.

\textbf{However}, at the redshifts probed by DESI ($z \sim 0.1$--2.1),
the transition regime $g_N/\aM \sim 1$--10 applies, softening the
enhancement. The predicted DFD/\LCDM{} ratio for $f\sigma_8$ is:
\begin{equation}
  \frac{(f\sigma_8)_{\rm DFD}}{(f\sigma_8)_{\Lambda\rm CDM}}
  = 1.05\text{--}1.15,
\end{equation}
which is detectable at $\sim 2$--$3\sigma$ across 4 combined DESI
redshift bins.

\textbf{Pre-registration target:} Submit a DFD prediction paper
with explicit $f\sigma_8(z_i)$ predictions for each DESI bin by
August 2026, before Euclid DR1 cosmology release (October 2026).

\subsection{BAO Amplitude Prediction (P27)}

Prediction P27 (introduced at i23): the BAO peak amplitude in the
matter power spectrum scales as:
\begin{equation}
  A_{\rm BAO} \propto \sqrt{\aM}
\end{equation}
in the DFD framework, due to the $\sqrt{\aM}$ enhancement in the
growth factor. This is a \emph{unique} DFD signature that is absent
in \LCDM{} and standard MOND theories without cosmological perturbation
theory.

\textbf{Status:} Not yet testable without mochi\_class. Analytic
estimates suggest $A_{\rm BAO}^{\rm DFD}/A_{\rm BAO}^{\Lambda\rm CDM}
\approx 1.03$--1.10, which would be detectable with DESI DR2 precision.

\subsection{Neutrino Mass Constraints}

DESI DR2 + CMB yields the tightest cosmological bound:
\begin{equation}
  \sum m_\nu < 0.064 \text{ eV} \quad (95\% \text{ CL, } \Lambda\text{CDM}).
\end{equation}
In $w_0w_a$CDM: $\sum m_\nu < 0.16$ eV.

KATRIN 2025 direct measurement: $m_{\nu_e} < 0.45$ eV (90\% CL),
improved from $< 0.8$ eV.

\textbf{DFD implication:} DFD's enhanced growth rate means that a
given $\sigma_8$ can be achieved with \emph{lower} neutrino masses
than in \LCDM{} (since enhanced growth compensates for neutrino
free-streaming suppression). DFD predicts:
\begin{equation}
  \sum m_\nu^{\rm DFD} \approx 0.06\text{--}0.10 \text{ eV},
\end{equation}
consistent with both cosmological and direct bounds.

\subsection{Euclid Timeline}

Euclid Q1 data release (March 2025): 63 deg$^2$ of calibrated images
and catalogs, focused on ``legacy science'' (not cosmology). As of
March 2025, Euclid has observed $\sim 2000$ deg$^2$ ($\sim 14\%$ of
total survey area).

\textbf{Cosmology data release: October 2026.} This will include
the first weak lensing $S_8$ measurement and growth rate constraints
from Euclid. DFD's $f\sigma_8$ predictions must be pre-registered
before this date.

\subsection{New Literature Reviewed}

\begin{tcolorbox}[colback=green!5!white, colframe=green!50!black,
  title=\textbf{Agent 12: NEW LITERATURE REVIEWED}]

\begin{enumerate}

\item \textbf{Li et al.\ (2026), ``Status of the $S_8$ Tension:
  A 2026 Review of Probe Discrepancies,'' arXiv:2602.12238.}
  Comprehensive review finding S$_8$ tension is ``probe-dependent.''
  KiDS-Legacy consistent with Planck; DES Y6 remains $2.7\sigma$ low.
  \emph{Supporting:} DFD's $\sigma_8 \approx 0.83$ consistent with
  converging probes.

\item \textbf{Wright et al.\ (2025), ``KiDS-Legacy: Cosmological
  constraints from cosmic shear,'' A\&A 2025.}
  KiDS-Legacy $S_8 = 0.815 \pm 0.016$, dramatic upward shift from
  KiDS-1000. \emph{Supporting:} eliminates previous DFD tension
  with KiDS.

\item \textbf{DES Collaboration (2025), ``DES Y6 cosmic shear results.''}
  DES Y6 $S_8$ remains low at $2.4$--$2.7\sigma$ below Planck.
  \emph{Mild tension:} DFD's $S_8 \approx 0.83$ is $\sim 2.5\sigma$
  above DES Y6.

\item \textbf{KATRIN Collaboration (2025), ``Direct neutrino-mass
  measurement based on 259 days of KATRIN data,'' Science.}
  $m_{\nu_e} < 0.45$ eV (90\% CL). \emph{Neutral:} consistent with
  DFD's prediction of $\sum m_\nu \sim 0.06$--0.10 eV.

\item \textbf{DESI Collaboration (2025), ``DESI DR2 BAO + CMB
  neutrino mass,'' arXiv:2503.14738.}
  $\sum m_\nu < 0.064$ eV (95\% CL) in \LCDM.
  \emph{Consistent:} with DFD's $0.06$--0.10 eV prediction range.

\end{enumerate}
\end{tcolorbox}

%=============================================================================
\section{Corrections to Previous Iterations}
%=============================================================================

\begin{tcolorbox}[colback=yellow!5!white, colframe=yellow!70!black,
  title=\textbf{CORRECTIONS FROM i24}]

\begin{enumerate}

\item \textbf{i23 CMB deficit range corrected:} i23 stated 15--30\% deficit
  as the best case. After careful review of the transition-regime
  enhancement factor (1.6$\times$ vs.\ deep-MOND 40--140$\times$),
  the full uncertainty range is \textbf{10--61\%}, not 15--30\%. The
  lower bound requires all four mechanisms operating at their maximum,
  which is optimistic but not excluded.

\item \textbf{i23 Euclid DR1 $f\sigma_8$ ratio:} i23 quoted 1.05--1.15.
  This is confirmed but should be understood as the prediction for
  \emph{combined} bins, not individual bins where the ratio may be
  outside this range.

\item \textbf{No other corrections needed.} All other i23 findings
  remain valid.

\end{enumerate}
\end{tcolorbox}

%=============================================================================
\section{Summary and Next Steps}
%=============================================================================

\begin{enumerate}
  \item \textbf{mochi\_class remains the \#1 priority.} DESI DR2 and
    Euclid DR1 cosmology (October 2026) create an urgent timeline.
  \item \textbf{$\psi$-screening as dark energy:} DESI DR2's $3.1\sigma$
    preference for dynamical dark energy is a significant opportunity
    for DFD. Quantitative predictions require $\Sigma(y)$ resolution.
  \item \textbf{Pre-register $f\sigma_8$ predictions} for DESI bins
    and Euclid by August 2026.
  \item \textbf{Cosmic birefringence:} DFD's $\beta = 0$ prediction
    strengthened by Planck/ACT miscalibration discrepancy.
  \item \textbf{S$_8$ tension dissolving:} DFD's prediction sits
    comfortably in the converging measurement landscape.
\end{enumerate}

\end{document}
